\documentclass[../main.tex]{subfiles}
\begin{document}

\section{Bối cảnh và động lực}
\label{section:1.1}

Các ứng dụng mô phỏng lấy bối cảnh đô thị Việt Nam cần mô phỏng được hành vi đặc trưng của người tham gia giao thông tại Việt Nam. Một mô hình tái hiện được hành vi đó tại nút giao còn cho phép đánh giá bằng mô phỏng những phương án mà hiện nay chỉ đánh giá được sau khi nút giao đã xây dựng hoặc biện pháp đã được áp dụng trên thực tế. Mô phỏng giao thông vi mô tái hiện từng phương tiện và tương tác giữa chúng trên một mạng đường chi tiết tới từng làn và từng hướng rẽ~\cite{lopez2018sumo}, nên một thiết kế hay một biện pháp có thể được thử nghiệm trước khi triển khai. Trong mô phỏng giao thông đô thị, mô hình hành vi như vậy phục vụ ít nhất ba nhóm ứng dụng.

Nhóm thứ nhất là quy hoạch đô thị và thiết kế nút giao. Khi quy hoạch một nút giao mới, chưa có dòng xe thật để quan sát; người thiết kế chỉ có thể so sánh các thiết kế hình học, cách chia làn và chương trình đèn bằng hàng chờ, độ trễ và số xung đột mà mô phỏng tạo ra. Nếu mô hình giả định mọi người dừng đúng vạch và rẽ đúng làn, phương án được chọn chỉ tốt cho một dòng xe không tồn tại trên thực tế. Nhóm thứ hai là tổ chức giao thông: các biện pháp như cấm một hướng rẽ, thay đổi pha đèn hay bổ sung làn rẽ phải liên tục ở một nút đang khai thác cần được ước lượng tác động trước khi áp dụng. Nhóm thứ ba là kiểm thử hệ thống hỗ trợ lái và xe tự hành. Semrau và Erdmann xây một khung mô phỏng để thử hệ thống hỗ trợ lái trong các tình huống giao thông đặc thù của Trung Quốc~\cite{semrau2016chinese}, vì một hệ thống chỉ được kiểm thử với dòng xe tuân thủ thì chưa được kiểm thử với những tình huống nó sẽ gặp trên thực tế. Một xe tự hành đi qua nút giao ở Hà Nội sẽ gặp xe máy vượt đèn đỏ theo đoàn, xe máy xuất phát trước đèn xanh và xe máy đứng chờ giữa nút; các tình huống đó cần được tái hiện trong mô phỏng để có thể kiểm thử.

Kết quả của các ứng dụng trên chỉ đáng tin cậy khi mô hình hành vi sát với thực tế. Một mô phỏng cho hàng chờ ngắn hơn, ít xung đột hơn hoặc ít tình huống nguy hiểm hơn thực tế sẽ đánh giá quá cao thiết kế đang thử. Luận văn xây dựng một phương pháp để có được mô hình hành vi như vậy cho một nút giao có đèn: từ một danh mục hành vi có cơ sở trong tài liệu, một mô hình ngôn ngữ lớn (LLM) đề xuất bộ luật hành vi, mã tất định thực thi bộ luật trong mô phỏng, và sai lệch giữa mô phỏng và dữ liệu quan sát được dùng để sửa bộ luật. Sản phẩm của luận văn là một khung công cụ mã nguồn mở kèm một bộ luật mẫu đã hiệu chỉnh, dùng lại được cho các nhóm ứng dụng trên.

\section{Phát biểu bài toán}
\label{section:1.current}

Các trình mô phỏng vi mô phổ biến dùng mô hình hành vi chuẩn: xe bám theo xe phía trước trong cùng làn, dừng trước vạch khi đèn đỏ, chỉ vào nút khi được phép và nhường đường theo quyền ưu tiên. Mô hình làn con cho phép nhiều xe hẹp chiếm chung một làn và di chuyển theo phương ngang, nên tái hiện được phần không gian của dòng xe máy~\cite{semrau2016chinese}. Các mô hình này được xây dựng cho dòng xe đi theo làn và tuân thủ luật. Với tham số mặc định, chúng không tạo ra vượt đèn đỏ có xác suất, vượt đèn theo đoàn, xuất phát sớm, hay xe rẽ vào hướng cấm rồi dừng chờ giữa nút.

Giao thông Việt Nam khác dòng xe đó ở hai điểm. Điểm thứ nhất là mật độ và thành phần dòng xe. Xe máy chiếm trên 90\% số phương tiện đăng ký tại Hà Nội~\cite{nguyen2022motorcycles} và khoảng 86,5\% số phương tiện đếm được tại một nút có đèn trong giờ cao điểm~\cite{mai2019signal}; ô tô và xe máy dùng chung lòng đường, và xe máy chen vào khoảng trống giữa các phương tiện khác thay vì đi theo làn cố định. Điểm thứ hai là mức tuân thủ. Video tại hai nút ở Hà Nội cho thấy 9,64--11,06\% xe máy vượt đèn đỏ và 13,55--15,82\% xe máy xuất phát từ khoảng 5~s trở lên trước khi đèn xanh, trong khi các tỷ lệ này ở ô tô chỉ từ 0,50\% đến 1,61\%~\cite{mai2019signal}. Cùng nghiên cứu ghi nhận nhiều xe máy của luồng phụ đi vào giữa nút và dừng chờ khi pha đèn chuyển.

Theo hiểu biết của tác giả, trong phạm vi khảo sát ở Chương~\ref{chapter:Related_works}, chưa có nghiên cứu nào mô phỏng các hành vi đặc trưng của người tham gia giao thông Việt Nam tại một nút giao có đèn. Các nghiên cứu mô phỏng nút giao hiện có dùng mô hình chuẩn, nên không bao gồm các hành vi không tuân thủ nêu trên. Công trình gần nhất là mô hình dựa trên tác tử của Saval và cộng sự cho ô tô và xe máy ở Hà Nội~\cite{saval2023mixed}. Mô hình này đưa được hành vi không tuân thủ vào mô phỏng bằng luật xác suất, nhưng cấu trúc luật do người lập mô hình chọn cố định, tham số được hiệu chỉnh cho một đoạn phố thay vì cho một nút, và mô hình không có xuất phát sớm, vượt đèn theo đoàn hay rẽ trái dừng chờ giữa nút. Các hướng hiệu chỉnh tự động, như tối ưu hộp đen~\cite{akiba2019optuna}, tìm được giá trị số cho tham số nhưng cần cấu trúc luật có sẵn. Các công trình dùng LLM với trình mô phỏng~\cite{jeong2025agentsumo,li2026chatsumo} hướng tới cấu hình kịch bản và nhu cầu, không sinh luật hành vi tài xế.

Bài toán của luận văn là mô phỏng hành vi giao thông tại một nút giao, phát biểu như sau. Cho một nút giao có đèn với dòng xe hỗn hợp ô tô và xe máy, mật độ cao và có hành vi không tuân thủ, trong một khung giờ (giờ cao điểm hoặc giờ thấp điểm), cần xây dựng một bộ luật hành vi người tham gia giao thông sao cho trình mô phỏng chạy với bộ luật đó tái hiện các hành vi quan sát được tại nút trong khung giờ đó với sai lệch nhỏ nhất. Đầu vào gồm mô tả nút giao (hình học, làn, chương trình đèn, hướng cấm) và dữ liệu quan sát (số đếm theo hướng, sự kiện vi phạm, hàng chờ, vị trí xe). Đầu ra là một bộ luật kiểm tra được bằng máy, thực thi được trong mô phỏng và tái lập được với cùng seed.

Bài toán có ba khó khăn. Thứ nhất, cấu trúc của mô hình hành vi chưa biết trước: ngoài giá trị của một tập tham số cố định, cần xác định luật nào được dùng, điều kiện kích hoạt của từng luật và nhóm người lái mà luật áp dụng. Thứ hai, không gian tham số lớn, các tham số tác động lẫn nhau, và mỗi lần đánh giá phải chạy mô phỏng nhiều lần với nhiều seed. Thứ ba, số liệu đo về hành vi tại nút giao ở Việt Nam còn rời rạc; nhiều đại lượng cần cho hiệu chỉnh chưa được công bố.

\section{Định hướng giải pháp}
\label{section:1.3}

Giải pháp của luận văn tách việc \emph{sinh luật} khỏi việc \emph{thực thi luật} (Hình~\ref{fig:1.1}).

Sinh luật là việc xác định thành phần của mô hình hành vi: chia người tham gia giao thông thành các nguyên mẫu tài xế với tỷ trọng và tham số riêng, chọn các luật có điều kiện (ví dụ ``vượt đèn đỏ với xác suất tăng theo số người vừa vượt'') và giá trị tham số của chúng. Việc này do một LLM chạy cục bộ thực hiện, ngoài vòng mô phỏng. LLM nhận mô tả nút giao, danh mục hành vi và loại luật, danh mục tham số được phép chỉnh kèm miền giá trị, và sai lệch giữa mô phỏng và dữ liệu quan sát của vòng trước. LLM trả về một đặc tả luật dạng văn bản có cấu trúc, kèm lý do cho từng lựa chọn.

Thực thi luật là việc điều khiển hành vi của các xe trong mô phỏng theo đặc tả. Việc này do mã tất định thực hiện, không có sự tham gia của LLM. Đặc tả được kiểm tra hợp lệ, sau đó được biên dịch thành tham số tĩnh mà trình mô phỏng nạp khi khởi động và luật có điều kiện mà một bộ điều khiển đánh giá ở mỗi bước mô phỏng. Bộ điều khiển chỉ can thiệp trong các giới hạn an toàn định trước và ghi lại mọi lần can thiệp. Kết quả mô phỏng được chuyển về cùng dạng với dữ liệu quan sát để tính chỉ số; các chỉ số trở thành phản hồi cho lần sinh luật kế tiếp.

\begin{figure}[!htbp]
\centering
\begin{tikzpicture}[
    node distance=8mm and 10mm,
    box/.style={draw, rounded corners=2pt, align=center, minimum height=11mm, text width=2.7cm, font=\small},
    llm/.style={box, fill=gray!20, very thick},
    arr/.style={-{Stealth[length=2.4mm]}, thick}
]
\node[llm] (llm) {Sinh luật:\\LLM cục bộ};
\node[box, right=of llm] (val) {1. Kiểm tra\\hợp lệ};
\node[box, right=of val] (cmp) {2. Biên dịch};
\node[box, below=of cmp] (sim) {3. Mô phỏng và\\bộ thực thi luật};
\node[box, left=of sim] (met) {4. Tính chỉ số\\M1--M10};
\node[box, below=of llm] (fb) {5. Phản hồi\\cho vòng sau};
\draw[arr] (llm) -- (val);
\draw[arr] (val) -- (cmp);
\draw[arr] (cmp) -- (sim);
\draw[arr] (sim) -- (met);
\draw[arr] (met) -- (fb);
\draw[arr] (fb) -- (llm);
\begin{pgfonlayer}{background}
\node[draw, dashed, gray, rounded corners=3pt, inner sep=4mm, fit=(val)(cmp)(sim)(met), label={[font=\scriptsize, gray]above:thực thi luật: mã tất định, không gọi LLM}] {};
\end{pgfonlayer}
\end{tikzpicture}
\caption{Một vòng hiệu chỉnh. LLM chỉ chạy ở bước sinh luật (Chương~\ref{chapter:RuleGeneration}); các bước 1--4 là mã tất định (Chương~\ref{chapter:RuleExecution}). Bước 2 tạo ra tham số tĩnh của các nguyên mẫu tài xế và tập luật có điều kiện cho bộ thực thi.}
\label{fig:1.1}
\end{figure}

Cách tách hai phần này phù hợp với giao thông Việt Nam hơn cách tiếp cận thông thường ở các nước có dòng xe tuân thủ, vì bốn lý do.

Lý do thứ nhất liên quan đến phần thay đổi của mô hình giữa các kịch bản. Với dòng xe tuân thủ, hành vi tại nút giao do luật giao thông và một số mô hình đã được nghiên cứu kỹ quyết định: bám xe, chuyển làn, chấp nhận khoảng trống. Cấu trúc mô hình giống nhau ở mọi nút; hiệu chỉnh là chỉnh giá trị tham số của cấu trúc đó cho khớp số đo tại nút~\cite{hollander2008calibration}. Với giao thông Việt Nam, phần quyết định độ giống thực tế là các hành vi lệch khỏi luật, và hành vi nào xảy ra phụ thuộc vào nút và khung giờ: nút có vỉa hè đi được hay không, có hướng cấm rẽ hay không, độ dài hàng chờ giờ cao điểm. Cấu trúc luật vì vậy phải được chọn lại cho từng kịch bản, và đây là phần cần hiểu biết về hành vi người lái. Luận văn đặt giả thuyết rằng LLM có thể đề xuất phần này, còn bộ tối ưu số thì không, vì bộ tối ưu số cần cấu trúc có sẵn. Ngược lại, phần vật lý của chuyển động (bám xe, tránh va chạm, cơ chế nhường đường) do mô hình chuẩn của trình mô phỏng đảm nhiệm: tham số của nó được hiệu chỉnh, nhưng cấu trúc của nó không thay đổi giữa các kịch bản. Tách hai phần cho phép giao phần cấu trúc thay đổi cho LLM và giữ phần cơ chế trong mã tất định.

Lý do thứ hai: hành vi vi phạm dễ gây ra va chạm hoặc hiện tượng xe đi xuyên qua nhau trong mô phỏng, vì nó buộc phải gỡ một số ràng buộc của mô hình chuẩn. Với dòng xe tuân thủ, vấn đề này hầu như không phát sinh. Khi việc thực thi được tách riêng, các giới hạn an toàn được tập trung tại một thành phần và áp dụng như nhau cho mọi luật, bất kể nguồn sinh luật.

Lý do thứ ba: mỗi bước mô phỏng của một nút giao giờ cao điểm có hàng trăm xe, phần lớn là xe máy. Gọi LLM cho từng xe ở từng bước không khả thi về thời gian và không tái lập được. Khi luật được sinh một lần và thực thi bằng mã tất định, mỗi lần chạy tái lập được với cùng seed.

Lý do thứ tư: đặc tả luật là văn bản có cấu trúc, nên kiểm tra được trước khi chạy, so sánh được giữa các phương án, đọc được bởi người làm quy hoạch, và về nguyên tắc chuyển được sang trình mô phỏng khác bằng một bộ biên dịch khác.

\section{Mục tiêu và phạm vi}
\label{section:1.2}

\subsection{Mục tiêu}

Mục tiêu của luận văn là xây dựng một phương pháp mô phỏng hành vi người tham gia giao thông Việt Nam tại nút giao có đèn, trong đó LLM sinh và tinh chỉnh lặp bộ luật hành vi, còn mã tất định thực thi bộ luật trong mô phỏng. Luận văn có năm mục tiêu cụ thể: (i) xác định các hành vi cần mô hình hóa và cơ sở của từng hành vi trong tài liệu; (ii) định nghĩa đặc tả luật gồm tham số tĩnh và luật có điều kiện, kèm bộ kiểm tra hợp lệ; (iii) xây dựng bộ thực thi luật chỉ can thiệp trong giới hạn an toàn và ghi lại mọi lần can thiệp; (iv) xây dựng hệ thống chỉ số để so sánh mô phỏng với dữ liệu quan sát và cập nhật luật; và (v) so sánh luật do LLM sinh với mô hình mặc định, bộ luật viết tay và một bộ tối ưu số trên cùng ngân sách mô phỏng.

\subsection{Phạm vi}

Phương pháp xét một nút giao có đèn với dòng xe gồm ô tô và xe máy, trong một khung giờ (giờ cao điểm hoặc giờ thấp điểm). Người đi bộ, xe buýt, xe tải, ô tô đỗ dưới lòng đường và việc chọn lộ trình không được mô hình hóa. Các hành vi được xét chia thành ba nhóm theo cơ chế: không tuân thủ tín hiệu, tìm khoảng trống khi ùn, và tương tác trong nút (Mục~\ref{section:rg_behaviours}). LLM sinh luật hành vi nhưng không sinh hay sửa mạng đường; các không gian chỉ xe vi phạm được dùng (ví dụ vỉa hè) là bản vá tĩnh của mạng.

Chương~\ref{chapter:RuleGeneration} và Chương~\ref{chapter:RuleExecution} trình bày phương pháp cho một nút giao bất kỳ, không gắn với địa điểm hay giờ cụ thể. Chương~\ref{chapter:Experiment} áp dụng phương pháp cho một nút giao cụ thể ở Hà Nội trong một khung giờ cao điểm để minh họa và đánh giá. Dữ liệu quan sát lấy từ hai đoạn video ngắn quay tại nút trong giờ cao điểm. Từ video đo được lưu lượng, tỷ lệ hướng đi, tỷ lệ xe máy và chu kỳ đèn; các chỉ số hành vi mà video chưa đo được lấy giá trị mục tiêu từ số liệu đã công bố về các nút giao khác, mỗi giá trị gắn nhãn chuyển giao hoặc giả định. Kết quả thực nghiệm vì vậy cho biết quy trình vận hành được với nhu cầu đo tại nút và chỉ ra những điểm mô phỏng khác mô hình chuẩn; mức độ sát thực tế của các hành vi vi phạm chỉ kiểm chứng được khi chính các hành vi đó được đo từ video. Khung giờ thấp điểm chưa được thực nghiệm.

\section{Câu hỏi nghiên cứu}
\label{section:rq}

Các mục tiêu trên dẫn tới ba câu hỏi nghiên cứu mà phần thực nghiệm phải trả lời. Việc đánh giá mức độ sát thực tế của mô phỏng đòi hỏi các chỉ số hành vi đo bằng video tại nút. Do luận văn mới đo được từ video lưu lượng, tỷ lệ hướng đi và chu kỳ đèn, câu hỏi đó nằm ngoài ba câu hỏi dưới đây.

\begin{enumerate}[label=\textbf{CH\arabic*.}, leftmargin=*]
    \item Bộ thực thi luật có tạo ra trong mô phỏng các hành vi vi phạm mà mô hình chuẩn của trình mô phỏng không có, mà không làm tăng va chạm mô phỏng và số xe kẹt so với mô hình chuẩn không?
    \item Trên cùng không gian tham số và cùng ngân sách số lần chạy mô phỏng, đặc tả do LLM cục bộ sinh (chỉ tham số, và cả cấu trúc lẫn tham số) so với đặc tả do bộ tối ưu số tìm ra như thế nào?
    \item Một LLM cục bộ quy mô chín tỷ tham số có sinh được đặc tả hợp lệ không, mắc những lỗi gì, và bộ kiểm tra hợp lệ có chặn được các lỗi đó trước khi tiêu tốn lần chạy mô phỏng không?
\end{enumerate}

\section{Đóng góp của luận văn}
\label{section:1.contrib}

Luận văn có bốn đóng góp chính. Bảng~\ref{table:contrib} nối mỗi đóng góp với câu hỏi nghiên cứu, bằng chứng và mức kiểm chứng hiện có.
\begin{enumerate}
    \item Một phương pháp mô phỏng hành vi tại nút giao tách sinh luật khỏi thực thi luật, trong đó LLM chỉ sinh đặc tả luật ngoài vòng mô phỏng và mọi đặc tả đi qua cùng một bộ kiểm tra hợp lệ và một bộ thực thi.
    \item Một danh mục 11 hành vi của người tham gia giao thông tại nút có đèn, xây dựng bằng cách tổng hợp các nghiên cứu đã công bố và bổ sung từ quan sát, kèm cơ chế kích hoạt và biến điều khiển của luật tương ứng.
    \item Các nguyên tắc thực thi luật an toàn và kiểm chứng được: chỉ gỡ đúng ràng buộc mà hành vi vi phạm cần gỡ, chỉ giữ xe, không đẩy xe qua xung đột, và ghi lại việc trình mô phỏng có thi hành từng lệnh hay không.
    \item Một hệ thống chỉ số và khung so sánh các nguồn sinh đặc tả trên cùng không gian tham số và cùng ngân sách mô phỏng, đã hiện thực trên SUMO (Simulation of Urban MObility)~\cite{lopez2018sumo} và chạy với LLM cục bộ \texttt{qwen3.5:9b} trên nhu cầu đo từ video tại nút.
\end{enumerate}

\begin{table}[!htbp]
\centering
\small
\caption{Đóng góp, câu hỏi nghiên cứu, bằng chứng và mức kiểm chứng.}
\label{table:contrib}
\begin{tabular}{|p{1.4cm}|p{1.3cm}|>{\raggedright\arraybackslash}p{5.3cm}|>{\raggedright\arraybackslash}p{5.1cm}|}
\hline
\textbf{Đóng góp} & \textbf{Câu hỏi} & \textbf{Bằng chứng} & \textbf{Mức kiểm chứng} \\ \hline
1 & CH2, CH3 & Quy trình vận hành trọn vòng với Optuna và LLM cục bộ (Mục~\ref{subsection:campaigns_n1108}, \ref{subsection:qwen_campaign}) & Vận hành được với nhu cầu đo từ video; chưa có đặc tả nào đạt tiêu chí chấp nhận \\ \hline
2 & -- & Tổng hợp tài liệu (Mục~\ref{section:rg_behaviours}); mười hành vi hiện thực trên SUMO (Bảng~\ref{table:3.1}) & Có cơ sở tài liệu cho hành vi 1--4, 9, 10; hành vi 5--8, 11 từ quan sát, chưa đo \\ \hline
3 & CH1 & Thí nghiệm speed mode của SUMO, kiểm thử giữ xe, nhật ký can thiệp (Mục~\ref{section:5.2}, \ref{section:5.3}) & Kiểm chứng bằng thí nghiệm và kiểm thử tự động \\ \hline
4 & CH1, CH2 & So sánh B0 và B1 (Mục~\ref{section:compare_b0_b1}); các chiến dịch cùng ngân sách (Mục~\ref{subsection:campaigns_n1108}) & Khung so sánh tuân thủ ngân sách; phép so sánh LLM với Optuna chưa cho kết luận \\ \hline
\end{tabular}
\end{table}

Mục~\ref{section:solutions} trình bày bảy giải pháp kỹ thuật được dùng để hiện thực các đóng góp này trên SUMO.

\section{Bố cục luận văn}
\label{section:1.4}

Phần còn lại của luận văn được tổ chức như sau.

Chương~\ref{chapter:Related_works} trình bày nền tảng lý thuyết: các khái niệm cơ bản (pha và chu kỳ của đèn tín hiệu, thời gian sau xâm phạm, mô hình bám xe và chuyển làn, thống kê GEH), số liệu đã công bố về hành vi xe máy tại nút có đèn, mô phỏng vi mô với mô hình làn con, và các công trình về mô hình hóa dòng xe hỗn hợp, hiệu chỉnh mô phỏng, và LLM trong tối ưu hóa và mô phỏng giao thông.

Chương~\ref{chapter:RuleGeneration} nêu tổng quan phương pháp rồi trình bày phần sinh luật: định nghĩa kịch bản, danh mục 11 hành vi được mô hình hóa, cơ sở khoa học và biến điều khiển của từng loại luật, đặc tả luật, lý do và cách dùng LLM để sinh luật, và vòng tinh chỉnh theo phản hồi.

Chương~\ref{chapter:RuleExecution} trình bày phần thực thi luật: kiểm tra hợp lệ, hai lớp thực thi, không gian riêng cho xe vi phạm, các nguyên tắc can thiệp an toàn, nhật ký can thiệp, hệ thống chỉ số đo kết quả của một lần chạy, và thiết kế so sánh các nguồn sinh đặc tả.

Chương~\ref{chapter:Experiment} hiện thực phương pháp trên SUMO cho nút Lê Văn Lương -- Hoàng Minh Giám -- Nguyễn Tuân (Hà Nội) trong giờ cao điểm, mô tả dữ liệu quan sát lấy từ video quay tại nút và cấu hình đánh giá, rồi trình bày kết quả theo ba câu hỏi nghiên cứu và đánh giá các bộ luật tốt nhất trên ngày kiểm định.

Chương~\ref{chapter:conclusion} tóm tắt kết quả đạt được, trả lời từng câu hỏi nghiên cứu, nêu hạn chế và hướng phát triển.

\end{document}
